\section{Isolation for maximal symplectic holonomy}\label{sec:maximal}
Let \(C\) be a closed Riemann surface of genus \(g\ge2\). Consider a germ, along a marked horizontal section, of a flat fibration with symplectic four-dimensional fibers. Its local transition maps are symplectomorphisms independent of the base coordinate. On the horizontal distribution lift the complex structure of \(C\), and on the vertical distribution take any smooth compatible complex structure. The derivative of holonomy at the marked point gives a symplectic representation \(\rho:\pi_1(C)\to\mathrm{Sp}(4,\mathbb R)\) with associated flat bundle \(W\). Assume the Calabi--Yau condition along the section,
\[c_1(W)[C]=2g-2 .\]
In this normalization it is the positive maximal value for a surface-group symplectic representation of rank four. We prove that one neighborhood isolates the marked curve in \emph{all} multiplicities, that its full series depends only on the deformation class of the linear holonomy, and that different linear holonomies can already be distinguished by the oriented count of the underlying curve. This statement concerns the smooth universal cycle algebra \(S\).

\subsection{Uniform isolation and deformation to the linear representation}

There is no nonzero vector with bounded orbit under \(\rho\). Suppose there is such a vector. On its span under all translates the representation has compact closure: choose finitely many translates as a basis; every group element and its inverse are uniformly bounded on them. If the restricted symplectic form on this invariant subspace has a nonzero radical, the radical is a nonzero flat isotropic subbundle \(I\). The quotient \(I^\perp/I\) is a flat symplectic bundle of real rank at most two. Topologically the remaining symplectic summand is \(I\oplus I^\vee\) and its first Chern class vanishes on \(C\). With a metric it is the complexification of a real bundle, whose first Chern class is two-torsion; \(H^2(C;\mathbb Z)\) is torsion-free. Milnor's inequality bounds the absolute Euler number of the remaining rank-two flat bundle by \(g-1\). If the restricted form has zero radical, the bounded subrepresentation is itself symplectic with compact holonomy, hence admits a flat unitary metric and has first Chern class zero on \(C\). The symplectic orthogonal is again of rank at most two. Both alternatives contradict the displayed Chern number. The same proof applies to every finite-index subgroup by multiplicativity of the Chern number and of \(g-1\) \cite{milnor}.

Choose a norm on a transverse Darboux chart. For each unit vector \(v\) there is a holonomy word \(w\) with \(\lVert\rho(w)v\rVert>2\). The strict inequalities and compactness of the sphere give a finite set of such words covering the sphere. In a compact smooth parameter family of marked germs with this Chern number one may do the same over the entire unit-sphere bundle in finitely many parameter charts. All chosen words are defined on a common smaller transverse ball. Their nonlinear maps differ from their derivatives by \(o(\lVert x\rVert)\) uniformly. For every sufficiently small nonzero \(x\) at least one word therefore takes it to norm greater than \(3\lVert x\rVert/2\). A nonzero finite invariant set contained in the ball is impossible by applying this at its point of maximum norm. Shrinking the tube about the entire marked section ensures that any horizontal curve contained in it determines such a full finite invariant set. In fact the same finite-word argument rules out any nonzero full orbit trapped in a strictly smaller ball by iteration.

The closed form given by the fiber symplectic form in the flat charts vanishes on the marked section. It is exact on the tube since that tube retracts onto the section. It is nonnegative on complex lines and vanishes precisely on horizontal lines. Any compact irreducible holomorphic curve in the tube is therefore horizontal and must be the marked curve. For every finite or branched cover of \(C\), a solution of the pullback linear normal equation has a holomorphic graph in the flat vector bundle. Its vertical area vanishes by the radial primitive, so the solution is parallel. Its monodromy image has finite index by compactness and openness of the lifted covering map, hence there are no nonzero solutions. The original normal operator is the one from \(\rho\) with the chosen compatible structure along the section. Its Fredholm index on the identity cover is zero, so this marked map is regular.

In transverse Darboux charts conjugate every transition \(F\) by positive real dilations, using \(F_t(v)=t^{-1}F(tv)\) for \(0<t\le1\) and \(F_0=dF(0)\). These are symplectic local maps satisfying the same cocycle identities; they extend smoothly at zero on a common smaller tube. The pulled-back vertical complex structures likewise extend with value at the zero section. Equivalently these charts are the rescaling of the given smooth fibration germ along its marked section. The finite-word argument and the exact-area argument apply uniformly throughout the parameter interval. In every degree \(m\) the multiple section is therefore a compact open family of cycles, and its characteristic class is constant by the compact prism relation. Interpolating two compatible vertical structures gives the same result. A piecewise smooth path of actual maximal representations also gives a family of linear flat bundles and the identical compact prism. Conjugation changes the transversal by an isomorphism. Thus to each such deformation class one can assign the universal series of the marked support. The statements about consistency along paths and loops here concern a \emph{persistent marked leaf} with this degree condition. Birth and loss of horizontal sections in an unmarked affine deformation are different questions. The proof requires no uniform genus bound on moduli spaces of stable maps.

For the double of a Fuchsian spin lift the value in every multiplicity is \(E_g\) by the elementary-neighborhood theorem, as in Section~\ref{sec:horizontal}. More generally, if the two identical summands are mixed by an oriented flat rotation representation on their multiplicity space, one may deform the rotation character of \(\pi_1(C)\) to the trivial character in its torus of characters. The Chern number is unchanged along this path and the entire series is again \(E_g\). The component decomposition of maximal symplectic representations is nontrivial; the preceding arguments do not identify the series for every component \cite{gothen} \cite{gm}.

\subsection{An oriented count distinguishing two maximal models}

Take the sum of two spin lifts of the same uniformizing representation, with theta characteristics \(\theta_1,\theta_2\) on \(C\). The Chern number, underlying smooth symplectic normal-bundle type and geometric support all agree with those of a doubled model. The determinant orientation of the regular core map is
\[(-1)^{h^0(\theta_1)+h^0(\theta_2)} .\]

For one lift let \(z\in\mathbb H\), \(y=\operatorname{Im} z\), and write \((a,b)\) for real flat fiber coordinates. Give the plane its positive-degree complex structure using the complex coordinate \(s=(a-\bar z b)/y\). Under a matrix \(\left(\begin{smallmatrix}A&B\\c&d\end{smallmatrix}\right)\) this transforms by \(s\mapsto(cz+d)s\), so \(s^2dz\) is an ordinary differential and \(s\) is a spin coordinate. In this frame the normal graph operator and an interpolation to the holomorphic spin operator are
\[D_t s=\bar\partial s+ t\,\frac{i\bar s}{2y}\,d\bar z,
                           \qquad 0\le t\le1 .\]
For \(t=1\) the formula follows by differentiating the coordinate of a constant flat vector and using the Leibniz rule for a real Cauchy--Riemann operator. If \(D_t s=0\), then \(d(s^2dz)\) is \(2t|s|^2\,dx\wedge dy/y\) upstairs, hence defines a nonnegative density on \(C\). Stokes forces \(s=0\) for \(t>0\). The real Fredholm index is zero, so all these operators are invertible.

At zero the kernel is \(H^0(\theta)\) and the cokernel is \(H^1(\theta)=H^0(\theta)^\vee\). The first derivative between these spaces pairs \(s\) with \(w\) by integration of \(w\bar s/y\) times a nonzero constant against oriented coordinate area. This is the antilinear isomorphism of a positive Hermitian form up to that constant. In the canonical complex orientations on kernel and cokernel its real determinant has sign \((-1)^{h^0(\theta)}\). In a kernel-complement splitting the other Fredholm block is invertible and has the complex orientation at zero; the reduced map is \(t\) times this isomorphism plus \(O(t^2)\). This computes the canonical orientation sign at small positive \(t\), and invertibility up to \(t=1\) preserves it. Signs multiply for the two summands. Changing auxiliary compatible complex structures in the flat linear model preserves invertibility by the exact-area argument.

Even and odd theta characteristics both exist in every genus \(g\ge2\). Choosing opposite parities gives a negative regular mark; choosing the same lift twice gives a positive one. Their degree-one characteristic insertions are \emph{different elements of \(S\)}. If they agreed, proper pullback at source genus \(g\) would leave just their two regular embeddings, by the minimum genus and the unstable-leaf argument. The compact witness construction of \cite{WueProduct} transports these points to constant parameter balls in one family. Its oriented local functional gives their respective determinant signs, contradicting the asserted equality. The regular map has trivial stabilizer in both cases. This distinguishes the full local series without assuming a universal numerical evaluation, and without identifying the negative series with an inverse elementary series. The familiar spin parity enters through an orientation of an actual stable map \cite{WueProduct} \cite{atiyah}.

For opposite parities the quotient \(\theta_2/\theta_1\) is a nontrivial flat two-torsion line. On its connected unramified double cover, the two spin lifts become identical. The entire pulled-back series is therefore \(E_{2g-1}\). The original negative count and this positive elementary pullback give a concrete case in which the covering operation carries information beyond the scalar count on the base. The finite-cover signs and the resulting least-genus coefficients are computed in Appendix~\ref{sec:flat-sign}.
